% Pregunta 3 (2004-junio-OEP2004): dos rectas de balance y preferencia revelada.
% Redibujado del escaneo del examen. Autor: Víctor Gutiérrez Marcos
\documentclass[tikz,border=8pt]{standalone}
\usepackage[T1]{fontenc}
\usepackage{newpxtext,newpxmath}
\begin{document}
\begin{tikzpicture}[x=1.3cm,y=1.3cm,line cap=round]
  \draw[->,>=latex] (0,0) -- (6.4,0) node[below=4pt] {$x_1$};
  \draw[->,>=latex] (0,0) -- (0,4.3) node[left=4pt] {$x_2$};
  \draw[thick] (0,3.6) -- (1.32,0);   % RB1: muy inclinada
  \draw[thick] (0,2.2) -- (5.1,0);    % RB2: más plana
  \fill (0.3,2.78) circle (1.6pt);
  \node[right] at (0.55,2.9) {$(x_1, x_2)$};
  \fill (4.3,0.345) circle (1.6pt);
  \node[above right] at (4.25,0.4) {$(y_1, y_2)$};
  \fill (2.05,0.95) circle (1.6pt);
  \node[below] at (2.05,0.85) {$(z_1, z_2)$};
  \node (rb1) at (0.75,4.15) {$RB_1$};
  \draw[->,>=latex] (rb1.south west) -- (0.23,3.0);
  \node (rb2) at (2.6,2.3) {$RB_2$};
  \draw[->,>=latex] (rb2.south) -- (2.25,1.3);
\end{tikzpicture}
\end{document}
